\subsection{自同态族的分解}

设 V 是一个向量空间，S 是一个集合，r 是从 S 到 End(V) 中的一个映射。用 P 表示从 S 到 k 的映射所成的集合。若 \(\lambda \in P\) ，则用 \(\mathrm{V}_{\lambda}(\mathrm{S})\) （相应地，\(\mathrm{V}^{\lambda}(\mathrm{S})\) ）表示满足下列条件的 \(v \in V\) 所成的集合：对所有 \(s \in S\) ，\(r(s)v = \lambda(s)v\) （相应地，\((r(s) - \lambda(s))^{n}v = 0\) 对充分大的 n 成立）。集合 \(\mathrm{V}_{\lambda}(\mathrm{S})\) 与 \(\mathrm{V}^{\lambda}(\mathrm{S})\) 都是 V 的向量子空间，且 \(\mathrm{V}_{\lambda}(\mathrm{S}) \subset \mathrm{V}^{\lambda}(\mathrm{S})\) 。我们称 \(\mathrm{V}_{\lambda}(\mathrm{S})\) 为 V 相对于 \(\lambda\) （以及 r）的特征子空间，称 \(\mathrm{V}^{\lambda}(\mathrm{S})\) 为 V 相对于 \(\lambda\) （以及 r）的准素子空间，并称 \(\mathrm{V}^{0}(\mathrm{S})\) 为 V 的（相对于 r 的）幂零空间。我们称 \(\lambda\) 是 S 在 V 中的一个权，如果 \(\mathrm{V}^{\lambda}(\mathrm{S}) \neq 0\) 。

特别地，若 S 只含单个元素 s，则 P 可与 k 等同；这时我们用记号 \(\mathrm{V}_{\lambda(s)}(s)\) 与 \(\mathrm{V}^{\lambda(s)}(s)\) ，或 \(\mathrm{V}_{\lambda(s)}(r(s))\) 与 \(\mathrm{V}^{\lambda(s)}(r(s))\) ，来代替 \(\mathrm{V}_{\lambda}(\{s\})\) 、\(\mathrm{V}^{\lambda}(\{s\})\) ；我们说 \(r(s)\) 的特征子空间、准素子空间与幂零空间；\(\mathrm{V}_{\lambda(s)}(s)\) 中的元素 v 称为 \(r(s)\) 的特征向量，且当 \(v \neq 0\) 时，\(\lambda(s)\) 称为相应的特征值（参看~代数，第七章，§5）。

对所有 \(\lambda\in P\) ，下列关系式是直接的：

\[
V ^ {\lambda} (S) = \bigcap_ {s \in S} V ^ {\lambda (s)} (s),\tag{1}
\]

\[
V _ {\lambda} (S) = \bigcap_ {s \in S} V _ {\lambda (s)} (s).\tag{2}
\]

设 \(k'\) 是 \(k\) 的一个扩张。从 \(\operatorname{End}(V)\) 到 \(\operatorname{End}(V \otimes_k k')\) 的典则映射与 \(r\) 复合，给出一个映射 \(r' : S \to \operatorname{End}(V \otimes_k k')\) 。类似地，任何映射 \(\lambda\) （从 S 到 \(k\) ）都典则地定义一个映射，仍记作 \(\lambda\) （从 S 到 \(k'\) ）。在这些记号下，我们有以下命题：

命题 1. 对所有 \(\lambda \in \mathrm{P}\) ，

\[
(\mathrm{V} \otimes_ {k} k ^ {\prime}) ^ {\lambda} (\mathrm{S}) = \mathrm{V} ^ {\lambda} (\mathrm{S}) \otimes_ {k} k ^ {\prime} \text {and} (\mathrm{V} \otimes_ {k} k ^ {\prime}) _ {\lambda} (\mathrm{S}) = \mathrm{V} _ {\lambda} (\mathrm{S}) \otimes_ {k} k ^ {\prime}.
\]

设 \((a_{i})\) 是 k-向量空间 \(k'\) 的一个基。若 \(v \in V \otimes_{k} k'\) ，则 v 可唯一地表示为 \(\sum v_{i} \otimes a_{i}\) 的形式，其中 \((v_{i})\) 是 V 中元素的一个有限支集族。对所有 \(s \in S\) ，

\[
(r ^ {\prime} (s) - \lambda (s)) ^ {n} (v) = \sum (r (s) - \lambda (s)) ^ {n} v _ {i} \otimes a _ {i}.
\]

由此得出

\[
v \in (\mathrm{V} \otimes_ {k} k ^ {\prime}) ^ {\lambda} (\mathrm{S}) \Longleftrightarrow v _ {i} \in \mathrm{V} ^ {\lambda} (\mathrm{S}) \text { for   all } i,
\]

\[
v \in (\mathrm{V} \otimes_ {k} k ^ {\prime}) _ {\lambda} (\mathrm{S}) \Longleftrightarrow v _ {i} \in \mathrm{V} _ {\lambda} (\mathrm{S}) \text { for   all } i,
\]

这就蕴含该命题。

命题 2. 设 \(\mathrm{V}, \mathrm{V}'\) 、\(\mathrm{W}\) 为向量空间。设 \(r: \mathrm{S} \to \operatorname{End}(\mathrm{V})\) 、\(r': \mathrm{S} \to \operatorname{End}(\mathrm{V}')\) 与 \(q: \mathrm{S} \to \operatorname{End}(\mathrm{W})\) 为映射。

\begin{enumerate}
\def\labelenumi{(\roman{enumi})}
\item
  设 \(f: \mathrm{V} \to \mathrm{W}\) 是线性映射，使得 \(q(s)f(v) = f(r(s)v)\) 对 \(s \in \mathrm{S}\) 与 \(v \in \mathrm{V}\) 成立。那么，对所有 \(\lambda \in \mathrm{P}\) ，\(f\) 把 \(\mathrm{V}^{\lambda}(\mathrm{S})\) （相应地，\(\mathrm{V}_{\lambda}(\mathrm{S})\) ）映入 \(W^{\lambda}(\mathrm{S})\) （相应地，\(W_{\lambda}(\mathrm{S})\) ）。
\item
  设 \(\mathrm{B}:\mathrm{V}\times \mathrm{V}'\to \mathrm{W}\) 是双线性映射，使得
\end{enumerate}

\[
q (s) \mathrm{B} (v, v ^ {\prime}) = \mathrm{B} (r (s) v, v ^ {\prime}) + \mathrm{B} (v, r ^ {\prime} (s) v ^ {\prime})
\]

其中 \(s \in \mathrm{S}\) ，\(v \in \mathrm{V}\) ，\(v' \in \mathrm{V}'\) 。那么，对所有 \(\lambda, \mu \in \mathrm{P}\) ，B 把 \(V^{\lambda}(\mathrm{S}) \times \mathrm{V}'^{\mu}(\mathrm{S})\) （相应地，\(\mathrm{V}_{\lambda}(\mathrm{S}) \times \mathrm{V}_{\mu}'(\mathrm{S})\) ）映入 \(W^{\lambda + \mu}(\mathrm{S})\) （相应地，\(W_{\lambda + \mu}(\mathrm{S})\) ）。

\begin{enumerate}
\def\labelenumi{(\roman{enumi})}
\setcounter{enumi}{2}
\tightlist
\item
  设 \(B: V \times V' \to W\) 是双线性映射，使得
\end{enumerate}

\[
q (s) \mathrm{B} (v, v ^ {\prime}) = \mathrm{B} (r (s) v, r ^ {\prime} (s) v ^ {\prime})
\]

其中 \(s \in S\) ，\(v \in V\) ，\(v' \in V'\) 。那么，对所有 \(\lambda, \mu \in P\) ，B 把 \(V^{\lambda}(S) \times V'^{\mu}(S)\) （相应地，\(V_{\lambda}(S) \times V'_{\mu}(S)\) ）映入 \(W^{\lambda\mu}(S)\) （相应地，\(W_{\lambda\mu}(S)\) ）。

在情形 (i) 中，\((q(s) - \lambda(s))^n f(v) = f((r(s) - \lambda(s))^n v)\) 对 \(s \in S\) 与 \(v \in V\) 成立，由此即得结论。在情形 (ii) 中，

\[
(q (s) - \lambda (s) - \mu (s)) \mathrm{B} (v, v ^ {\prime}) = \mathrm{B} ((r (s) - \lambda (s)) v, v ^ {\prime}) + \mathrm{B} (v, (r ^ {\prime} (s) - \mu (s)) v ^ {\prime})
\]

其中 \(s \in S\) ，\(v \in V\) ，\(v' \in V'\) ，从而对 n 用归纳法得

\[
(q (s) - \lambda (s) - \mu (s)) ^ {n} \mathrm{B} (v, v ^ {\prime}) = \sum_ {i + j = n} \binom {n} {i} \mathrm{B} ((r (s) - \lambda (s)) ^ {i} v, (r ^ {\prime} (s) - \mu (s)) ^ {j} v ^ {\prime}).
\]

(ii) 中的断言立即可得。在情形 (iii) 中，

\[
(q (s) - \lambda (s) \mu (s)) \mathrm{B} (v, v ^ {\prime}) = \mathrm{B} ((r (s) - \lambda (s)) v, r ^ {\prime} (s) v ^ {\prime}) + \mathrm{B} (\lambda (s) v, (r ^ {\prime} (s) - \mu (s)) v ^ {\prime})
\]

其中 \(s \in S\) ，\(v \in V\) ，\(v' \in V'\) ，从而对 n 用归纳法得

\[
\begin{array}{l} (q (s) - \lambda (s) \mu (s)) ^ {n} \mathrm{B} (v, v ^ {\prime}) \\ = \sum_ {i + j = n} \binom {n} {i} \mathrm{B} (\lambda (s) ^ {j} (r (s) - \lambda (s)) ^ {i} v, r ^ {\prime} (s) ^ {i} (r ^ {\prime} (s) - \mu (s)) ^ {j} v ^ {\prime}). \end{array}
\]

(iii) 中的断言立即可得。

命题 3. 和 \(\sum_{\lambda \in \mathrm{P}}\mathrm{V}^{\lambda}(\mathrm{S})\) 是直和。和 \(\sum_{\lambda \in \mathrm{P}}\mathrm{V}_{\lambda}(\mathrm{S})\) 是直和。

第二个断言是第一个断言的推论；因此只需证明后者。我们区分几种情形。

\begin{enumerate}
\def\labelenumi{\alph{enumi})}
\item
  S 为空集。断言是平凡的。
\item
  S 只含单个元素 s。设 \(\lambda_{0}, \lambda_{1}, \ldots, \lambda_{n}\) 是 k 中互异的元素。对 \(i = 0, 1, \ldots, n\) ，设 \(v_{i} \in V^{\lambda_{i}}(s)\) 且假定 \(v_{0} = v_{1} + \cdots + v_{n}\) 。只需证明 \(v_{0} = 0\) 。对 \(i = 0, \ldots, n\) ，存在整数 \(q_{i} > 0\) 使得 \((r(s) - \lambda_{i})^{q_{i}} v_{i} = 0\) 。考虑多项式 \(\mathrm{P}(\mathrm{X}) = \prod_{i \geq 1} (\mathrm{X} - \lambda_{i})^{q_{i}}\) 与 \(\mathrm{Q}(\mathrm{X}) = (\mathrm{X} - \lambda_{0})^{q_{0}}\) 。我们有 \(\mathrm{Q}(r(s)) v_{0} = 0\) ，以及 \(\mathrm{P}(r(s)) v_{0} = \sum_{i=1}^{n} \mathrm{P}(r(s)) v_{i} = 0\) 。由于 P 与 Q 互素，Bezout 恒等式表明 \(v_{i} = 0\) 。
\end{enumerate}

\[
v _ {0} = 0
\]

\begin{enumerate}
\def\labelenumi{\alph{enumi})}
\setcounter{enumi}{2}
\item
  S 有限且非空。我们对 S 的基数作归纳。设 \(s \in \mathrm{S}\) ，并设 \(\mathrm{S}' = \mathrm{S} - \{s\}\) 。设 \((v_{\lambda})_{\lambda \in \mathrm{P}}\) 是 V 中元素的一个有限支集族，满足 \(\sum_{\lambda \in \mathrm{P}} v_{\lambda} = 0\) 且 \(v_{\lambda} \in \mathrm{V}^{\lambda}(\mathrm{S})\) 。取 \(\lambda_0 \in \mathrm{P}\) 。设 \(\mathrm{P}'\) 是 \(\lambda \in \mathrm{P}\) 中使 \(\lambda |\mathrm{S}' = \lambda_0|\mathrm{S}'\) 的元素所成的集合。对 \(\mathrm{S}'\) 应用归纳假设，得 \(\sum_{\lambda \in \mathrm{P}'} v_{\lambda} = 0\) 。若 \(\lambda, \mu\) 是 \(\mathrm{P}'\) 中互异的元素，则 \(\lambda(s) \neq \mu(s)\) 。由于和 \(\sum_{\alpha \in k} \mathrm{V}^{\alpha}(s)\) 依 \(b\) ) 是直和，且 \(v_{\lambda} \in \mathrm{V}^{\lambda(s)}(s)\) ，故 \(v_{\lambda} = 0\) 对所有 \(\lambda \in \mathrm{P}'\) 成立，特别地 \(v_{\lambda_0} = 0\) ，此即所要证明的。
\item
  一般情形。设 \((v_{\lambda})_{\lambda\in\mathrm{P}}\) 是 V 中元素的一个有限支集族，满足 \(\sum_{\lambda\in\mathrm{P}} v_{\lambda}=0\) 且 \(v_{\lambda}\in\mathrm{V}^{\lambda}(\mathrm{S})\) 。设 \(P'\) 是 \(\lambda\in P\) 中使 \(v_{\lambda}\neq0\) 的元素所成的有限集合，并设 \(S'\) 是 S 的一个有限子集，使得条件 \(\lambda\in P'\) 、\(\mu\in P'\) 、\(\lambda|S'= \mu|S'\) 蕴含 \(\lambda=\mu\) 。我们有 \(v_{\lambda}\in\mathrm{V}^{\lambda|S'}(S')\) ；应用 c)，即知 \(v_{\lambda}=0\) 对 \(\lambda\in P'\) 成立，证毕。
\end{enumerate}

回顾一下，若 \(x \in \operatorname{End}(\mathrm{V})\) ，我们用 ad \(x\) 表示映射 \(y \mapsto xy - yx = [x, y]\) ，它从 \(\operatorname{End}(\mathrm{V})\) 到其自身。

引理 1. 设 \(x, y \in \operatorname{End}(\mathbf{V})\) 。

\begin{enumerate}
\def\labelenumi{(\roman{enumi})}
\item
  假定 \(\mathrm{V}\) 是有限维的。那么 \(x\) 可三角化当且仅当 \(\mathrm{V} = \sum_{a\in k}\mathrm{V}^a (x)\) 。
\item
  若存在整数 \(n\) 使 \((\mathrm{ad}x)^n y = 0\) ，则每个 \(\mathrm{V}^a(x)\) 在 \(y\) 下稳定。
\item
  假定 \(\mathrm{V}\) 是有限维的。若 \(\mathrm{V} = \sum_{a \in k} \mathrm{V}^a(x)\) 且每个 \(\mathrm{V}^a(x)\) 在 \(y\) 下稳定，则存在整数 \(n\) 使 \((\operatorname{ad} x)^n y = 0\) 。
\end{enumerate}

(i) 由代数，第七章，§5，第 2 节，命题 3 得出。

设 \(\mathrm{E} = \operatorname{End}(\mathrm{V})\) 。设 B 是双线性映射 \((u, v) \mapsto u(v)\) ，从 \(\mathrm{E} \times \mathrm{V}\) 到 V。由 ad \(x\) 的定义，

\[
x (\mathrm{B} (u, v)) = B (u, x (v)) + \mathrm{B} ((\operatorname{ad} x) (u), v)
\]

其中 \(x \in \mathrm{E}\) ，\(u \in \mathrm{E}\) ，\(v \in \mathrm{V}\) 。让 \(x\) 作用在 \(\mathrm{E}\) 上，作用经由 \(\operatorname{ad} x\) 。由命题 2 (ii)，\(\mathrm{B}(\mathrm{E}^0(x), \mathrm{V}^a(x)) \subset \mathrm{V}^a(x)\) 对所有 \(a \in k\) 成立。若 \((\operatorname{ad} x)^n y = 0\) ，则 \(y \in \mathrm{E}^0(x)\) ，故 \(y(\mathrm{V}^a(x)) \subset \mathrm{V}^a(x)\) ，这就证明了 (ii)。

为证明 (iii)，可把 V 换成 \(V^{a}(x)\) ，并把 x（相应地，y）换成它在 \(V^{a}(x)\) 上的限制。把 x 换成 x - a 后，可假定 x 是幂零的。于是 \((\operatorname{ad} x)^{2\dim V-1} = 0\) （第一章，§4，第 2 节），这就证明了 (iii)。

注. 上面的论证表明，若 V 是有限维的且存在整数 n 使 \((\operatorname{ad} x)^{n}y = 0\) ，则 \((\operatorname{ad} x)^{2\dim V-1}y = 0\) 。

以后我们说映射 \(r : S \to \text{End}(V)\) 满足条件 (AC)（“拟交换”）是指：

(AC) 对 S 的每一对元素 \((s, s')\) ，存在整数 \(n\) 使得

\[
(\operatorname{ad} r (s)) ^ {n} r \left(s ^ {\prime}\right) = 0.
\]

定理 1. 假定 V 是有限维的。下列条件等价：

\begin{enumerate}
\def\labelenumi{(\roman{enumi})}
\item
  条件 (AC) 满足，且对所有 \(s \in \mathrm{S}\) ，\(r(s)\) 可三角化。
\item
  对所有 \(\lambda \in \mathrm{P}\) ，\(\mathrm{V}^{\lambda}(\mathrm{S})\) 在 \(r(\mathrm{S})\) 下稳定，且 \(\mathrm{V} = \sum_{\lambda \in \mathrm{P}} \mathrm{V}^{\lambda}(\mathrm{S})\) 。
\end{enumerate}

若 \(V = \sum_{\lambda \in P} V^{\lambda}(S)\) ，则 \(V = \sum_{a \in k} V^{a}(s)\) 对所有 \(s \in S\) 成立，于是由引理 1 可知 (ii) 蕴含 (i)。假定条件 (i) 满足。引理 1 与公式 (1) 表明每个 \(V^{\lambda}(S)\) 在 \(r(S)\) 下稳定。剩下要证 \(V = \sum_{\lambda \in P} V^{\lambda}(S)\) 。我们对 dim V 作归纳。我们区分两种情形。

\(a)\) 对所有 \(s \in \mathrm{S}\) ，\(r(s)\) 只有一个特征值 \(\lambda(s)\) 。于是 \(\mathrm{V} = \mathrm{V}^{\lambda}(\mathrm{S})\) 。

\begin{enumerate}
\def\labelenumi{\alph{enumi})}
\setcounter{enumi}{1}
\tightlist
\item
  存在 \(s \in S\) 使 \(r(s)\) 至少有两个互异的特征值。于是 V 是诸 \(\mathrm{V}^{a}(s)\) （\(a \in k\) ）的直和，且对所有 a 有 \(\dim \mathrm{V}^{a}(s) < \dim \mathrm{V}\) 。每个 \(\mathrm{V}^{a}(s)\) 在 \(r(\mathrm{S})\) 下稳定，只需应用归纳假设即可。
\end{enumerate}

推论 1. 假定 V 是有限维的且条件 (AC) 满足。设 \(k'\) 是 \(k\) 的一个扩张。假定对所有 \(s \in S\) ，\(r(s) \otimes 1\) 作为 \(V \otimes_k k'\) 的自同态可三角化。设 \(P'\) 是从 \(S\) 到 \(k'\) 的映射所成的集合。那么 \(V \otimes_k k' = \sum_{\lambda' \in P'} (V \otimes_k k')^{\lambda'}(S)\) 。

设 \(r': S \to \operatorname{End}(V \otimes_k k')\) 是由 \(r\) 定义的映射。若 \(s_1, s_2 \in S\) ，则存在整数 \(n\) 使 \((\operatorname{ad} r(s_1))^n r(s_2) = 0\) ，从而 \((\operatorname{ad} r'(s_1))^n r'(s_2) = 0\) 。现在只需应用定理 1。

推论 2. 假定 V 是有限维的且条件 (AC) 满足。用 \(\mathrm{V}^{+}(\mathrm{S})\) 表示向量子空间 \(\sum_{s\in \mathrm{S}}\left(\bigcap_{i\geq 1}r(s)^{i}\mathrm{V}\right)\) 。那么：

\begin{enumerate}
\def\labelenumi{(\roman{enumi})}
\item
  \(\mathrm{V}^0 (\mathrm{S})\) 与 \(\mathrm{V^{+}(S)}\) 在 \(r(\mathrm{S})\) 下稳定；
\item
  \(\mathrm{V} = \mathrm{V}^{0}(\mathrm{S}) \oplus \mathrm{V}^{+}(\mathrm{S});\)
\item
  V 的每个在 \(r(\mathrm{S})\) 下稳定且满足 \(\mathrm{W}^0 (\mathrm{S}) = 0\) 的向量子空间 W 都包含于 \(\mathrm{V}^{+}(\mathrm{S})\) ；
\item
  \(\sum_{s\in S}r(s)\mathrm{V}^{+}(\mathrm{S}) = \mathrm{V}^{+}(\mathrm{S}).\)
\end{enumerate}

此外，\(V^{+}(S)\) 是 V 中唯一具有性质 (i) 与 (ii) 的向量子空间。对 k 的任何扩张 \(k'\) ，有 \((\mathrm{V} \otimes_{k} k')^{+}(\mathrm{S}) = \mathrm{V}^{+}(\mathrm{S}) \otimes_{k} k'\) 。

最后一个断言是直接的。于是，顾及命题 1，在证明其余断言时可假定 k 是代数闭的。由定理 1，\(V = \sum_{\lambda \in P} V^{\lambda}(S)\) ，且诸 \(V^{\lambda}(S)\) 在 \(r(S)\) 下稳定。若 \(s \in S\) ，则 \(r(s)|V^{\lambda}(S)\) 的特征多项式为 \((X - \lambda(s))^{\dim V^{\lambda}(S)}\) ；由此得出，\(\bigcap_{i \geq 1} r(s)^{i} V^{\lambda}(s)\) 为零，如果 \(\lambda(s) = 0\) ；而等于 \(V^{\lambda}(S)\) ，如果 \(\lambda(s) \neq 0\) ；因此，

\[
\mathrm{V} ^ {+} (\mathrm{S}) = \sum_ {\lambda \in \mathrm{P}, \lambda \neq 0} \mathrm{V} ^ {\lambda} (\mathrm{S}),\tag{3}
\]

这就证明了 (i)、(ii) 与 (iv)。若 W 是 V 的一个在 \(r(S)\) 下稳定的向量子空间，则 \(W = \sum_{\lambda \in P} W^{\lambda}(S)\) 且 \(W^{\lambda}(S) = W \cap V^{\lambda}(S)\) 。若 \(W^{0}(S) = 0\) ，则可见 \(W \subset V^{+}(S)\) ，这就证明了 (iii)。

设 \(V'\) 是 \(V\) 的一个在 \(r(S)\) 下稳定且满足 \(V' \cap V^0(S) = 0\) 的向量子空间。于是 \(V'^0(S) = 0\) ，故由 (iii) 得 \(V' \subset V^+(S)\) 。若此外还有 \(V = V^0(S) + V'\) ，则可见 \(V' = V^+(S)\) 。证毕。

我们有时把 \((\mathrm{V}^{0}(\mathrm{S}),\mathrm{V}^{+}(\mathrm{S}))\) 称为 V 的、或映射 \(r:S\to\operatorname{End}(\mathrm{V})\) 的 Fitting 分解。若 S 只含单个元素 s，我们用 \(\mathrm{V}^{+}(s)\) 或 \(\mathrm{V}^{+}(r(s))\) 代替 \(\mathrm{V}^{+}(\{s\})\) 。我们有 \(\mathrm{V}=\mathrm{V}^{0}(s)\oplus\mathrm{V}^{+}(s)\) ，\(\mathrm{V}^{0}(s)\) 与 \(\mathrm{V}^{+}(s)\) 在 \(r(s)\) 下稳定，\(r(s)|\mathrm{V}^{0}(s)\) 幂零而 \(r(s)|\mathrm{V}^{+}(s)\) 双射。

推论 3. 设 V 与 V' 是有限维向量空间，\(r: S \to \text{End}(V)\) 与 \(r': S \to \text{End}(V')\) 是满足条件 (AC) 的映射。设 \(f: V \to V'\) 是满射线性映射，使得 \(f(r(s)v) = r'(s)f(v)\) 对 \(s \in S\) 与 \(v \in V\) 成立。那么 \(f(V^{\lambda}(S)) = V'^{\lambda}(S)\) 对所有 \(\lambda \in P\) 成立。

鉴于命题 1，可化归到 \(k\) 是代数闭域的情形。由定理 1，我们有 \(\mathrm{V} = \bigoplus_{\lambda \in \mathrm{P}} \mathrm{V}^{\lambda}(\mathrm{S})\) ，\(\mathrm{V}' = \bigoplus_{\lambda \in \mathrm{P}} \mathrm{V}'^{\lambda}(\mathrm{S})\) ，以及 \(\mathrm{V}' = f(\mathrm{V}) = \sum_{\lambda \in \mathrm{P}} f(\mathrm{V}^{\lambda}(\mathrm{S}))\) 。最后，由命题 2 (i) 有 \(f(\mathrm{V}^{\lambda}(\mathrm{S})) \subset \mathrm{V}'^{\lambda}(\mathrm{S})\) ，由此即得推论。

命题 4. 假定 \(k\) 是完全域。设 \(V\) 是有限维向量空间，\(u\) 是 \(\operatorname{End}(V)\) 中的一个元素，\(u_s, u_n\) 是 \(u\) 的半单分量与幂零分量（代数，第七章，§5，第 8 节）。

\begin{enumerate}
\def\labelenumi{(\roman{enumi})}
\item
  对所有 \(\lambda \in k\) ，\(V^{\lambda}(u) = V^{\lambda}(u_s) = V_{\lambda}(u_s)\) 。
\item
  若 \(\mathrm{V}\) 具有一个代数结构且 \(u\) 是 \(\mathrm{V}\) 的导子，则 \(u_{s}\) 与 \(u_{n}\) 是 \(\mathrm{V}\) 的导子。
\item
  若 \(\mathrm{V}\) 具有一个代数结构且 \(u\) 是 \(\mathrm{V}\) 的自同构，则 \(u_{s}\) 与 \(1 + u_{s}^{-1}u_{n}\) 是 \(\mathrm{V}\) 的自同构。
\end{enumerate}

鉴于命题 1，可假定 \(k\) 是代数闭的，于是

\[
\mathrm{V} = \sum_ {\lambda \in k} \mathrm{V} ^ {\lambda} (u).
\]

\(u|V^{\lambda}(u)\) 的半单分量是比率为 \(\lambda\) 的、\(V^{\lambda}(u)\) 中的位似。这就证明了 (i)。

以下假定 \(\mathrm{V}\) 具有一个代数结构。设 \(x \in \mathrm{V}^{\lambda}(u)\) ，\(y \in \mathrm{V}^{\mu}(u)\) 。

若 \(u\) 是 \(\mathrm{V}\) 的导子，则 \(xy \in \mathrm{V}^{\lambda + \mu}(u)\) （命题 2 (ii)），于是

\[
u _ {s} (x y) = (\lambda + \mu) (x y) = (\lambda x) y + x (\mu y) = (u _ {s} x) y + x (u _ {s} y).
\]

这证明 \(u_{s}\) 是 V 的导子。于是 \(u_{n} = u - u_{s}\) 是 V 的导子。

若 u 是 V 的自同构，则 \(\operatorname{Ker}(u_{s}) = \mathrm{V}^{0}(u) = 0\) ，故 \(u_{s}\) 双射。另一方面，\(xy \in \mathrm{V}^{\lambda\mu}(u)\) （命题 2 (iii)），于是

\[
u _ {s} (x y) = (\lambda \mu) (x y) = (\lambda x) (\mu y) = u _ {s} (x). u _ {s} (y).
\]

这证明 \(u_{s}\) 是 V 的自同构；而如下元素也是自同构：

\[
1 + u _ {s} ^ {- 1} u _ {n} = u _ {s} ^ {- 1} u.
\]

